% ===== 原书 p.6 =====
\section{SYZ 变换}

由 Liouville--Arnold 定理，若有 Lagrange 环面纤维化 $\pi_Y:Y\to B$，则在其正则值集合 $B^{\mathrm{sm}}$ 上，局部地存在称为作用坐标（action coordinate）的坐标，除热带仿射变换（$\mathrm{GL}_n(\R)\ltimes\R^n$ 的作用称为实仿射变换，$\mathrm{GL}_n(\Z)\ltimes\R^n$ 的作用称为热带仿射变换，$\mathrm{GL}_n(\Z)\ltimes\Z^n$ 的作用称为整仿射变换）之外唯一确定。与这些作用坐标关于 $Y$ 的辛结构共轭的坐标，便是纤维的角坐标（angle coordinate）。这些局部坐标经热带仿射变换粘贴，在 $B^{\mathrm{sm}}$ 上定出热带仿射结构。$B^{\mathrm{sm}}$ 的切丛 $TB^{\mathrm{sm}}$ 与余切丛 $T^*B^{\mathrm{sm}}$ 上自然定出 $n$ 阶整局部系统 $T_{\Z}B^{\mathrm{sm}}$ 与 $T^*_{\Z}B^{\mathrm{sm}}$，逐纤维取商，便得到互为对偶的环面纤维丛 $X(B^{\mathrm{sm}}):=T^*B^{\mathrm{sm}}/T^*_{\Z}B^{\mathrm{sm}}$ 与 $\check X(B^{\mathrm{sm}}):=TB^{\mathrm{sm}}/T_{\Z}B^{\mathrm{sm}}$。$X(B^{\mathrm{sm}})$ 具有自然的辛结构，若 $\pi_Y$ 有 Lagrange 截面，则它与 $\pi_Y^{-1}(B^{\mathrm{sm}})$ 辛同胚。另一方面，$\check X(B^{\mathrm{sm}})$ 具有自然的复结构，但要把它（部分）紧化为 $B$ 上具有环面纤维化的复流形，就必须在这个自然复结构上加上量子修正。这个量子修正由对以 $\pi_Y$ 的 Lagrange 纤维为边界的全纯圆盘作计数得到，并期待能用 $B$ 的热带几何的语言来描述，但其细节仍在发展之中。

代替考虑 $\check X(B^{\mathrm{sm}})$ 复结构的量子修正，[AAK] 按如下方式构造了 $Y$ 的镜像流形 $\check Y$：首先，考虑从 $B$ 中除去 $\pi_Y\left(p^{-1}(U_Z\times\C)\right)$ 所得的子集 $B^{\mathrm{reg}}$。这里，$U_Z$ 是使 $U_Z\times\C$ 包含 $\chi$ 的支集 $U$ 的 $Z$ 的管状邻域。适当地取 $Z$ 的定义方程，并把 $U_Z$ 取得足够小，便可以使 $B^{\mathrm{reg}}$ 的连通分支与 $\Delta\cap\overline{M}$ 的元素一一对应。对于与各 $\alpha\in\Delta\cap\overline{M}$ 对应的 $B^{\mathrm{reg}}$ 的连通分支 $U_\alpha$，准备环面 $\widetilde U_\alpha:=\Spec\C\left[z_{\alpha,1}^{\pm1},\dots,z_{\alpha,n-1}^{\pm1},z_{\alpha,n}^{\pm1}\right]$，对 $\alpha,\beta\in\Delta\cap\Z^{n-1}$ 把 $\widetilde U_\alpha$ 与 $\widetilde U_\beta$ 按
\begin{equation}
\begin{gathered}
z_{\alpha,1}=(1+z_{\beta,n})^{\beta_1-\alpha_1}z_{\beta,1},\\
\vdots\\
z_{\alpha,n-1}=(1+z_{\beta,n})^{\beta_{n-1}-\alpha_{n-1}}z_{\beta,n-1},\\
z_{\alpha,n}=z_{\beta,n}.
\end{gathered}
\tag{13}
\end{equation}
\noindent 粘贴，再作适当的部分紧化，便得到 $\check Y$。(13) 由全纯圆盘的计数得到，是 Fukaya 范畴中的穿墙公式。
% [Tutor Note: 原文把粘贴指标集印作 $\alpha,\beta\in\Delta\cap\Z^{n-1}$，而同段前面用的是 $\Delta\cap\overline{M}$；$\Z^{n-1}$ 无法给出与 $z_{\alpha,n}$ 相配的 $\alpha_n,\beta_n$，(13) 末行更说明指标应属 $\overline{M}$。照录]

由同调镜像对称，期待 $X(B^{\mathrm{sm}})\to B^{\mathrm{sm}}$ 的 Lagrange 截面与 $\check X(B^{\mathrm{sm}})$ 上的线丛有如下关系：把 $B^{\mathrm{sm}}$ 的热带仿射坐标记作 $(x_i)_{i=1}^{n}$，则局部系统 $T_{\Z}B^{\mathrm{sm}}$ 与 $T^*_{\Z}B^{\mathrm{sm}}$ 分别由 $\{\partial/\partial x_i\}_{i=1}^{n}$ 与 $\{dx_i\}_{i=1}^{n}$ 生成。把对应于 $(x_i)_{i=1}^{n}$ 的 $TB^{\mathrm{sm}}$ 的纤维坐标记作 $(y_i)_{i=1}^{n}$，则关于 $\check X(B^{\mathrm{sm}})$ 的自然复结构，
\begin{equation}
z_i=\exp 2\pi\left(x_i+\sqrt{-1}y_i\right),\qquad i=1,\dots,n
\tag{14}
\end{equation}
\noindent 成为复坐标。把 $X(B^{\mathrm{sm}})\to B^{\mathrm{sm}}$ 的截面 $s$ 提升到 $T^*B^{\mathrm{sm}}\to B^{\mathrm{sm}}$，所得便定出 $B^{\mathrm{sm}}$ 上的一次微分形式 $\xi=\sum_{i=1}^{n}\xi_i(x)dx_i$。

% ===== 原书 p.7 =====
此时，在 $\check X(B^{\mathrm{sm}})$ 上的平凡线丛上引入由
\begin{equation}
\nabla:=d+2\pi\sqrt{-1}\sum_{i=1}^{n}\xi_i(x)dy_i
\tag{15}
\end{equation}
\noindent 给出的联络，所得的 $\nabla$ 称为 $s$ 的 SYZ 变换（SYZ transform）。$\nabla$ 的曲率为
\begin{equation}
F_\nabla=2\pi\sqrt{-1}\sum_{i,j}\frac{\partial\xi_j}{\partial x_i}dx_i\wedge dy_j
\tag{16}
\end{equation}
\noindent 其 $(2,0)$ 成分为
\begin{equation}
F^{(2,0)}_\nabla=\pi\sum_{i,j}\left(\frac{\partial\xi_j}{\partial x_i}-\frac{\partial\xi_i}{\partial x_j}\right)dz_i\wedge dz_j
\tag{17}
\end{equation}
\noindent 因此，「$s$ 是 Lagrange 截面」与「$\nabla$ 定出全纯结构」二者等价。把联络 $\nabla$ 限制到 $\check X(B^{\mathrm{sm}})\to B^{\mathrm{sm}}$ 在 $x\in B^{\mathrm{sm}}$ 上的纤维 $\check L_x$ 上，所得为平坦联络；这个平坦联络成为平凡联络的充要条件，是 $X(B^{\mathrm{sm}})\to B^{\mathrm{sm}}$ 的零截面 $s_0$ 在 $x\in B^{\mathrm{sm}}$ 处与 $s$ 相交。特别地，若 $\nabla$ 定出的 $\check X(B^{\mathrm{sm}})$ 上的线丛 $\mathcal L_\nabla$ 丰沛，且其第一 Chern 形式给出 $\check X(B^{\mathrm{sm}})$ 的辛形式，则 $\nabla$ 限制到 $\check L_x$ 上为平凡，正是 $\check L_x$ 为 Bohr--Sommerfeld 纤维的定义。这里，若存在范畴的等价 $D^b\mathrm{Fuk}\,X(B^{\mathrm{sm}})\cong D^b\mathrm{coh}\,\check X(B^{\mathrm{sm}})$，把零截面 $s_0$ 送到 $\mathcal O_{\check X(B^{\mathrm{sm}})}$、把截面 $s$ 送到 $\mathcal L_\nabla$，则
\begin{equation}
\mathrm{HF}^*(s_0,s)\cong\Hom^*\left(\mathcal O_{\check X(B^{\mathrm{sm}})},\mathcal L_\nabla\right)\cong H^*(\mathcal L_\nabla)
\tag{18}
\end{equation}
\noindent 特别地，若零截面 $s_0$ 与截面 $s$ 横截相交，且其 Maslov 指数全为 $0$，则 $s_0$ 与 $s$ 的 Floer 上同调 $\mathrm{HF}^*(s_0,s)$ 只含次数 $0$ 的部分，成为以 $s_0$ 与 $s$ 的交点为基的向量空间；而这些交点在上面的对应下被送到 $\check X(B^{\mathrm{sm}})\to B^{\mathrm{sm}}$ 的 Bohr--Sommerfeld 纤维。也就是说，$\check X(B^{\mathrm{sm}})$ 上的全纯线丛 $\mathcal L_\nabla$ 的全纯截面空间，具有与 $\check X(B^{\mathrm{sm}})\to B^{\mathrm{sm}}$ 的 Bohr--Sommerfeld 纤维相对应的基。注意这个断言中并未出现镜像 $X(B^{\mathrm{sm}})$。当 $B$ 紧、无奇异纤维、且 $\check X(B)$ 成为 Abel 流形时，这个基由 theta 函数给出。

要得到 $\check Y$ 上而非 $\check X(B^{\mathrm{sm}})$ 上的线丛，必须把 SYZ 变换的定义 (15) 修改得与穿墙公式 (13) 相容。作了适当的修改之后，可以证明下述结论：

\begin{theorem}[{[CPUa, Theorem 1.1]}]
对 $\check Y$ 上的任意线丛 $\mathcal F$，存在 $\pi_Y$ 的 Lagrange 截面 $L_{\mathcal F}$，使其 SYZ 变换为 $\mathcal F$，并且在 Hamilton 同痕意义下唯一。
\end{theorem}

证明的本质部分，是运用 [Abo09] 的技巧，从 $\check Y$ 的线丛出发，构造 $\overline{N}_{\C^\times}$ 的与 $Z$ 不相交的 Lagrange 子流形 $L$。由于 $Y$ 是 $\overline{N}_{\C^\times}\times\C$ 沿 $Z$ 的辛爆破的开子集，如前所述，把 $\R^{>0}$ 直积到 $L$ 上便可作出 $Y$ 的 Lagrange 子流形。此外，设定略有不同时的类似结果，是 [GM] 的主定理。

% CHUNK c03 DONE pages p6–p7
